\section{Idea 1: deletion turns a structural count into a recurrence}\label{sec:deletion}
\begin{definition}[Deletion outputs]\label{def:deletion}
For a domain \(D\) on \(X\), with \(n=|X|\), let
\[
 M=|D|,\quad D_{-x}=\{\sigma\text{ with }x\text{ erased}:\sigma\in D\},
 \quad M_x=|D_{-x}|,\quad\Sig(D)=\sum_{x\in X}M_x.
\]
Duplicates are removed in \(D_{-x}\). Outputs for different deleted
alternatives are counted separately, so an output is a pair \((x,P)\).
\end{definition}
\lean{\path{Condorcet.deleteEmbedding}, \path{deleteDomain},
\path{deletionSum}. Lean deletes from \(\mathrm{Fin}(k+1)\) and
canonically relabels the remaining alternatives as \(\mathrm{Fin}(k)\).
This relabelling changes neither distinctness nor any count.}

\begin{example}[Deletion collisions]\label{ex:four}
Take \(D=\{123,213,231,321\}\). It satisfies \(2N_3\), so it is
Condorcet and peak-pit. Its outputs are
\begin{center}
\begin{tabular}{c|ccc}\toprule
order & delete 1 & delete 2 & delete 3\\\midrule
123&23&13&12\\213&23&13&21\\231&23&31&21\\321&32&31&21\\\midrule
distinct outputs&2&2&2\\\bottomrule
\end{tabular}
\end{center}
Thus \(M=4\), \(\Sig=6\). Deleting 1 merges three orders into 23.
The useful structural theorems limit how much merging can happen on
average across all deletion outputs.
\end{example}
\outside{The displayed labels and small numerical example are expository
and are checked separately; no claim is made that this exact labelled
example is an exported Lean theorem. The general definitions and all
identities applied to it are formalised.}

\begin{lemma}[The deletion upper estimate]\label{lem:deletion-upper}
Restriction preserves both Condorcet and peak-pit domains. Consequently
\[
 \Sig(D)\leq nf(n-1)\quad\text{for Condorcet }D,\qquad
 \Sig(D)\leq nh(n-1)\quad\text{for peak-pit }D.
\]
\end{lemma}
\begin{proof}
Lift each order in any odd profile of restricted orders to one original
order of \(D\). All comparisons between surviving alternatives are
unchanged. Majority is therefore the restriction of a transitive
relation, hence transitive. An extreme triple law is also preserved by
restriction. Each of the \(n\) deletion families is thus bounded by
the relevant maximum at \(n-1\); sum those bounds.
\end{proof}
\lean{\path{Condorcet.IsCondorcet.restrict},
\path{IsPeakPit.restrict}, \path{deletionSum_le_general},
\path{deletionSum_le_peakPit}. The latter two literally use \(k+1\)
alternatives and a factor \(k+1\); set \(n=k+1\).}

\intuition{An upper bound on \(\Sig\) is easy because deletions have
one fewer alternative. The hard part is to prove a large \emph{lower}
bound on \(\Sig\) in terms of \(M\). If \(\Sig\geq c_nM\) for a
maximum domain, the two estimates give \(f(n)\leq(n/c_n)f(n-1)\).
The general algebraic observation is exposition; the final numerical
instance is an explicit Lean theorem.}

For perspective, \(M\leq\Sig\) holds for every order family, even
without the Condorcet property: count \(nM\) pairs \((\sigma,x)\),
and note that one output \((x,P)\) has at most \(n\) insertion positions
as preimages. This weak bound yields only a factorial recurrence.
The structural argument must force \(\Sig\) to be a positive fraction
of \(nM\).
